\documentclass[UTF8,a4paper,11pt]{ctexart}

% 建议使用 XeLaTeX 编译
\usepackage{geometry}
\usepackage{array}
\usepackage{tabularx}
\usepackage{longtable}
\usepackage{enumitem}
\usepackage{xcolor}
\usepackage{hyperref}
\usepackage{fancyhdr}
\usepackage{lastpage}

\geometry{left=2.4cm,right=2.4cm,top=2.3cm,bottom=2.3cm}
\setlength{\parindent}{0pt}
\setlength{\parskip}{0.45em}
\linespread{1.18}

\definecolor{mainblue}{RGB}{36,88,148}
\definecolor{lightgray}{RGB}{246,246,246}
\definecolor{midgray}{RGB}{120,120,120}

\hypersetup{
  colorlinks=true,
  linkcolor=mainblue,
  urlcolor=mainblue
}

\pagestyle{fancy}
\fancyhf{}
\lhead{\DocHeader}
\rhead{第 \thepage 页，共 \pageref{LastPage} 页}
\renewcommand{\headrulewidth}{0.4pt}

\setlist[enumerate]{leftmargin=2em,itemsep=0.35em,topsep=0.3em}
\renewcommand{\arraystretch}{1.45}
\setlength{\tabcolsep}{6pt}

\newcommand{\blankline}{\par\vspace{0.35em}\noindent\rule{\linewidth}{0.4pt}\par}
\newcommand{\blanklines}[1]{%
  \par
  \begingroup
  \newcount\i
  \i=0
  \loop\ifnum\i<#1
    \vspace{0.58em}\noindent\rule{\linewidth}{0.4pt}\par
    \advance\i by 1
  \repeat
  \endgroup
}
\newcommand{\fieldline}[1]{\textbf{#1}\quad\rule{0.82\linewidth}{0.4pt}\par}
\newcommand{\hint}[1]{{\small\color{midgray}#1}\par}
\newcommand{\sectionbar}[1]{%
  \vspace{0.7em}
  \noindent\colorbox{lightgray}{\parbox{0.97\linewidth}{\textbf{#1}}}
  \par\vspace{0.45em}
}
\newcommand{\basicinfo}[4]{%
\noindent\begin{tabularx}{\textwidth}{@{}lX lX lX lX@{}}
\textbf{#1} & \hrulefill & \textbf{#2} & \hrulefill & \textbf{#3} & \hrulefill & \textbf{#4} & \hrulefill
\end{tabularx}
\par\vspace{0.8em}
}
\newcolumntype{Y}{>{\raggedright\arraybackslash}X}
\newcolumntype{L}[1]{>{\raggedright\arraybackslash}p{#1}}

\newcommand{\DocHeader}{数值实验记录}

\title{\textbf{数值实验记录表}}
\author{姓名} % 自行修改
\date{日期} % 自行修改

\begin{document}
\maketitle

\section{基本信息}
\noindent\textbf{实验目的}\par \hint{一句话写明实验目的。}
\vspace{1.0em}
\noindent\textbf{主程序代码文件名}\par 
\vspace{1.0em}
\noindent\textbf{核心算法代码文件名}\par
\vspace{1.0em}
\noindent\textbf{其它}\par \hint{其它必要的说明}
\vspace{1.0em}



\section{基本设置}

\noindent\textbf{模型方程}
\hint{写清楚连续模型、优化问题或离散前的主要方程。重要公式可以直接写在下面。}
\vspace{1.0em}

\noindent\textbf{边界条件}
\hint{写清楚边界条件类型及其在程序中的处理方式。}
\vspace{1.0em}

\noindent\textbf{算法或数值格式}
\hint{例如梯度流、Newton 方法、半隐式格式、投影方法、有限差分格式、谱方法、有限元方法等。}
\vspace{1.0em}

\noindent\textbf{离散方法}
\hint{分别说明空间离散、时间离散、非线性项处理、约束或归一化处理，以及必要的后处理方法。}
\vspace{1.0em}

\noindent\textbf{其它}
\hint{记录与模型或算法有关但不方便归入以上几类的内容。}
\vspace{1.0em}


\section{参数设置}

\noindent\textbf{空间区域}\par
\hint{写清楚计算区域，例如 $[a,b]$、$[a,b]\times[c,d]$ 或三维区域。}
\vspace{1.0em}

\noindent\textbf{空间网格数或空间步长}\par
\hint{例如 $N=1000$，或 $\Delta x=0.01$。多维问题需分别写出各方向网格。}
\vspace{1.0em}

\noindent\textbf{时间步长}\par
\hint{例如 $\tau=10^{-4}$。若使用自适应时间步长，应写明取值范围和调整规则。}
\vspace{1.0em}

\noindent\textbf{终止条件}\par
\hint{例如终止时间、最大迭代步、残量阈值、能量变化阈值或其它停止准则。}
\vspace{1.0em}

\noindent\textbf{物理参数}\par
\hint{写清楚所有主要参数及其取值。}
\vspace{1.2em}

\noindent\textbf{初值}\par
\hint{写清楚初始函数、初始数据文件或初始扰动方式。}
\vspace{1.0em}

\noindent\textbf{边界条件}\par
\hint{如周期、Neumann、Dirichlet、自然边界条件、无通量条件等。若已在基本设置中写过，这里可只写实际程序采用的设置。}
\vspace{1.0em}

\noindent\textbf{其它}\par
\hint{包括随机种子、网格自适应参数、输出间隔、后处理参数等。}
\vspace{1.0em}

\section{实验结果与现象描述}
\hint{放置实验结果（图片/表格），并简述实验现象和结论。}

\vspace{2cm}


\section{下一步}
\hint{简述现有问题与下一步工作思路。如没有可以不写。}

\vspace{2cm}

\section{附录}
\hint{附上完整的代码文件包与必要的数据文件名，简述代码结构。}


\end{document}
